\documentclass[11pt,a4paper]{article}
\usepackage{fontspec}
\setmainfont{NotoSerifCJKkr-Regular.otf}[Path=fonts/,BoldFont=NotoSerifCJKkr-Bold.otf,ItalicFont=NotoSerifCJKkr-Regular.otf]
\setsansfont{NotoSerifCJKkr-Regular.otf}[Path=fonts/,BoldFont=NotoSerifCJKkr-Bold.otf]
\setmonofont{DejaVu Sans Mono}[Scale=0.79]
\XeTeXlinebreaklocale "ko"
\XeTeXlinebreakskip=0pt plus 1pt
\usepackage[margin=23mm,top=23mm,bottom=23mm,headheight=15pt]{geometry}
\usepackage{amsmath,amssymb,amsthm,mathtools}
\usepackage{booktabs,tabularx,longtable,array,graphicx,float}
\usepackage[dvipsnames]{xcolor}
\usepackage{enumitem,caption,fancyhdr,titlesec}
\usepackage{tikz}
\usetikzlibrary{arrows.meta,positioning,calc}
\usepackage{hyperref}
\hypersetup{unicode=true,colorlinks=true,linkcolor=MidnightBlue,urlcolor=MidnightBlue,citecolor=MidnightBlue,pdftitle={같은 COMMON의 Full--two-half 결함: 광자 chronology와 관측량 전달 정리},pdfauthor={WU088 HH 연구}}
\urlstyle{same}
\setlength{\emergencystretch}{3em}
\setlength{\parskip}{0.34em}
\setlength{\parindent}{1em}
\widowpenalty=10000
\clubpenalty=10000
\linespread{1.10}
\setlist{itemsep=0.18em,topsep=0.35em,leftmargin=1.6em}
\captionsetup{font=small,labelfont=bf}
\renewcommand{\figurename}{그림}
\renewcommand{\tablename}{표}
\renewcommand{\contentsname}{목차}
\renewcommand{\refname}{참고 문헌과 원자료}
\renewcommand{\proofname}{증명}
\titleformat{\section}{\large\sffamily\bfseries\color{MidnightBlue}}{\thesection}{0.7em}{}
\titleformat{\subsection}{\sffamily\bfseries}{\thesubsection}{0.7em}{}
\pagestyle{fancy}\fancyhf{}
\fancyhead[L]{\small\sffamily WU088\_HH / PHYS15 이론}
\fancyhead[R]{\small\sffamily 2026-10-13 KST}
\fancyfoot[C]{\thepage}
\newcommand{\dd}{\mathrm d}
\newcommand{\ex}{\operatorname{exact}}
\newcommand{\diag}{\operatorname{diag}}
\newcommand{\fh}{\widehat f}
\newcommand{\R}{\mathsf R}
\newcommand{\RF}{\mathsf R_F}
\newcommand{\BB}{\mathcal B}
\newcommand{\HF}{\mathcal H_F}
\newcommand{\nuH}{\nu_H}
\newcommand{\ew}{e_w}
\newcommand{\exx}{e_x}
\newcommand{\normi}[1]{\left\lVert #1\right\rVert_\infty}
\newtheorem{theorem}{정리}
\newtheorem{lemma}{보조정리}
\newtheorem{proposition}{명제}
\theoremstyle{definition}
\newtheorem{definition}{정의}
\newcommand{\ReviewDecision}{범위를 명시한 채택 (PROMOTE\_SCOPED)}
\newcommand{\ReviewExplanation}{독립 검토에서 9개 주장의 유도와 표적 검산을 확인하고 채택했다. 적용 범위는 원 captured 계수와 같은 COMMON을 공유하는 geometric isotropic source-aligned 계열이다. 실제 endpoint 수치 포함은 후속 계산 단위에 연결한다.}


\begin{document}
\begin{center}
{\color{MidnightBlue}\sffamily\large WU088\_HH \quad PHYS15 이론 연구}\\[5mm]
{\LARGE\bfseries 같은 COMMON의 Full--two-half 결함}\\[2mm]
{\Large\bfseries 광자 chronology와 관측량 전달 정리}\\[4mm]
{\large 2026년 10월 13일 KST}
\end{center}
\vspace{2mm}
\noindent\rule{\textwidth}{0.6pt}

\section*{초록}
같은 COMMON에서 출발하는 길이 $2h$의 Full과 실제 두 번의 길이 $h$ 단계를 결박하고, 두 endpoint의 차이를 공통 입력의 상관관계를 유지한 잔차로 구성한다. PHYS14의 등방 source-aligned 계열, 원 33개 에너지 node와 두 guards, 실제 두 번째 half의 근 포함을 재사용한다. 비교의 중심량은 실제 two-half 기체 $g_2$에서 평가한 Full residual $r=\HF(g_2)$다. 이 잔차를 source 시각, 중간 기체 이동, 일곱 광자 작용으로 정확히 분해한다.

광자 분해에는 opacity 변화, 흡수와 수송의 순서, 흡수 node의 시간 차이, BE 흡수 길이, remap 분할, 먼저 생성된 광자의 redshift 및 birth 흡수 시점이 들어간다. 원 cross section의 순서를 이용하면 중간 absorption--remap sector의 absorbed-photon 변화가 성분별로 비음수임을 보인다. 원 birth node의 추가 redshift sector는 이온화 absorption을 늘리면서 thermal deposition을 줄인다. 이 직접 forcing은 동일한 implicit inverse와 관측량 covector를 통해 endpoint 기여로 전달된다.

서로 다른 source 시각의 저장 밀도비를 보존하기 위해 온도를 일치시키는 gas 사상 $\Psi$를 구성한다. 원 FT03의 반응 인수 분해에서 nonphoto 및 HH source의 정확한 연결식을 얻는다. 기존 Full 전처리 행렬과 whole-domain 포함을 재사용하면 $\overline{\mathcal A}(g_2-g_F)=r$와 유효한 componentwise defect bound가 성립한다. 유한 온도 차이, 전자 수, 열에너지 및 photon endpoint는 이 같은 결함에서 계산된다. Number/energy guard 수지는 H 화학·열에너지 residual과 광자 chronology를 연결한다.

본 연구의 결과는 조건을 명시한 이론적 구성과 표적 대수 검산이다. 다음 계산 단위는 이미 확보된 source와 근 포함을 이 식에 대입하여 실제 $\Delta T,\Delta e,\Delta w,\Delta n_e$의 상관된 포함을 구하는 것이다.\cite{phys14,root_theory,photon_theory,bindings}

\vspace{1mm}
\noindent\fcolorbox{MidnightBlue!45}{MidnightBlue!3}{%
\parbox{\dimexpr\textwidth-2\fboxsep-2\fboxrule\relax}{%
\textbf{현재 연구 판정}\quad \ReviewDecision\par
\small \ReviewExplanation}}
\newpage
\begingroup
\small\setlength{\parskip}{0pt}
\tableofcontents
\vspace{3mm}
\noindent\textbf{읽기 순서.} 제1--3절에서 대상과 정확한 잔차를 정하고, 제4--6절에서 광자 기작을 분해한다. 제7절은 원 source 정규화, 제8--9절은 endpoint 포함과 관측량 전달을 다룬다. 제10--11절은 photon ledger와 실제 Full 광자 차이를 회수한다. 부록은 기호와 계산 입력의 대응을 모은다.
\endgroup
\newpage
\section{연구 대상과 재사용 전제}\label{sec:scope}
연구 질문은 같은 COMMON과 같은 미래 제어변수에서 Full과 실제 two-half의 기체·광자 endpoint 차이를 어떻게 구성하는가이다. PHYS14는 실제 첫 half의 기체 및 광자 carry를 둘째 half에 전달했다. 여기서는 그 완성된 경로와 같은 종료 시각의 Full을 비교한다. 핵심은 각 경로의 큰 포함을 독립적으로 뺄 때 반복되는 입력을 식에서 먼저 공유하는 것이다.

고정된 admissible 계수 선택을 $\alpha$라 하고, $\theta=(\lambda,b)\in[0,1]^2$를 두 경로에 동일하게 적용한다. $\lambda$는 HH 항, $b$는 광자 생성량에 곱한다. $\alpha$에는 COMMON, 과거 HH ON 이력, 원 source leaves, 기하학 및 정상화 angular weight가 포함된다. 모든 정리는 같은 $\alpha,\theta$에 대한 pointwise 식이다. 그 관계를 유지한 포함은 해당 계열 전체에 적용된다.

\begin{table}[H]
\centering\small
\begin{tabularx}{\textwidth}{@{}p{11mm}X p{49mm}@{}}
\toprule
전제 & 확인된 근거 & 이번 사용\\
\midrule
P1 & Full과 first half의 COMMON binary payload가 시간 간격 header를 제외하고 같다. & 같은 $g_0,Z_0,\lambda,b$\\
P2 & PHYS14의 geometric isotropic 계열과 정확한 39-state closure & $\R=\R(\rho)$, $\RF=\R(\rho^2)$\\
P3 & $g_1,g_2,g_F$가 원 convex source 영역 $X$ 내부에 있고 원 source 계열이 $C^2$다. & 연결 선분의 평균 Jacobian\\
P4 & Full과 둘째 half의 종료 시각·density·원자 leaves가 같다. & 공통 함수 $\BB_2$\\
P5 & Captured $M_F=2M$, HI active 16--32, He photo cross section은 0 & 같은 birth와 HI survival\\
P6 & Full의 고정 $C$와 whole-$X$ 전처리 derivative 포함 & 잔차에서 결함으로 전달\\
P7 & 원 source의 $\fh=n_{\rm He}^{\rm stored}/n_H$와 $C_T$ & source bridge와 관측량\\
\bottomrule
\end{tabularx}
\caption{PHYS09--14에서 재사용하는 전제와 PHYS15의 대응. 실제 파일 identity와 적용 범위는 원 binding 기록에 있다.}
\label{tab:premises}
\end{table}

시간과 단계 길이는
\begin{align}
t_0&=160000000000\ {\rm s},&
h&=625000000\ {\rm s},\label{eq:time}\\
t_1&=t_0+h,&
t_2&=t_0+2h=161250000000\ {\rm s}.
\end{align}
기하학은 세 방향의 captured $H_i=10^{-14}\,{\rm s}^{-1}$가 같은 등방 선택이다. 원 ray 표현에서 수송 비율은 $\rho=\exp[-\ex(H)\ex(h)]$로 공통이며, 실제 정상화 source weight의 합은 1이다. 이 선택이 기존 coefficient enclosure에 들어간다는 PHYS14의 geometry 검토를 재사용한다.\cite{phys14,bindings}

계산 상수는 원 IEEE binary64 leaf가 나타내는 정확한 실수로 해석한다. 표시용 소수와 underlying exact 값은 구별한다. 원 source 함수·근 포함과 이번 유도·표적 검산의 역할은 마지막 검증 절에서 연결한다.

\section{상태, 단위, source와 실제 단계식}\label{sec:stages}
\subsection{기체와 35차원 photon--guard 상태}
기체 좌표와 광자 상태를
\begin{equation}
g=(x,y_1,y_2,w)^{\mathsf T},\qquad
Z=(N_0,\ldots,N_{32},G_N,G_U)^{\mathsf T}
\label{eq:states}
\end{equation}
로 쓴다. $x$는 HII 분율, $y_1,y_2$는 HeII·HeIII 분율, $w$는 $\mathrm{eV/H}$다. $N_j,G_N$은 $\mathrm{photons/H}$, $G_U$는 $\mathrm{eV/H}$다. 전체 상태는 기체 4개와 photon--guard 35개, 합계 39개 실수로 닫힌다. 값과 $U=\partial_\lambda$, $V=\partial_b$, $W=\partial_\lambda\partial_b$의 네 carry slot에는 같은 대수 관계를 적용한다.

원 hat remap $\R$은 아래 경계로 나간 광자의 number와 energy를 두 guards에 더한다. 기존 guard energy는 $\rho$배 한다. 행벡터
\begin{equation}
\mathfrak n=(1,\ldots,1,1,0),\qquad
\mathfrak u=(E_0,\ldots,E_{32},0,1)
\label{eq:rows}
\end{equation}
는 각각 총 photon number와 energy를 읽는다. 따라서
\begin{equation}
\mathfrak n\R=\mathfrak n,\qquad \mathfrak u\R=\rho\mathfrak u.
\label{eq:remapmom}
\end{equation}
단위가 다른 photon--guard 성분의 norm이 필요하면 해당 단위와 척도를 함께 지정한다.

\subsection{생존·흡수 연산자와 birth}
Source 시각 $k=1,2$에서
\begin{equation}
a_{k,j}=\ex(c)\ex(n_k)\ex(\sigma_j),\qquad
\kappa_{k,j}(g)=a_{k,j}(1-x)
\label{eq:opacity}
\end{equation}
이며 단위는 $\mathrm{s}^{-1}$다. HI active group은 $j=16,\ldots,32$이다. Half와 Full의 survival 연산자는
\begin{align}
A_k(g)&=\diag\!\left[\frac1{1+h\kappa_{k,j}(g)}\right],
&L_k(g)&=I-A_k(g),\label{eq:halfA}\\
A_F(g)&=\diag\!\left[\frac1{1+2h\kappa_{2,j}(g)}\right],
&L_F(g)&=I-A_F(g).
\label{eq:fullA}
\end{align}
두 guard 성분에서는 $A=1,L=0$이다. 이 정의는 PHYS14의 source-aligned denominator를 따른다.

광흡수 한 개의 기체 방향은
\begin{equation}
\nu_j=(1,0,0,\epsilon_j)^{\mathsf T}.
\label{eq:nu}
\end{equation}
$\epsilon_j$는 원 captured HI heat leaf이다. $\mathsf V$의 active $j$열을 $\nu_j$로 두고 inactive 및 guard 열은 0으로 둔다.

Captured half birth amount와 Full amount는 정확히
\begin{equation}
M=\frac{7378697629483821}{2361183241434822606848},
\qquad M_F=2M,\qquad J=Me_{24}
\label{eq:birth}
\end{equation}
이다. 원 bits는 각각 \texttt{3eca36e2eb1c432d}, \texttt{3eda36e2eb1c432d}다. $M$의 단위는 photons/H이고 birth는 수송 뒤, 흡수 전에 들어간다.\cite{bindings}

\subsection{세 경로를 같은 COMMON에 결박한다}
실제 two-half의 absorption 전후 상태는
\begin{align}
Q_1&=\R Z_0+bJ,&
Z_1&=A_1(g_1)Q_1,&
P_1&=L_1(g_1)Q_1,\label{eq:stage1}\\
Q_2&=\R Z_1+bJ,&
Z_2&=A_2(g_2)Q_2,&
P_2&=L_2(g_2)Q_2.
\label{eq:stage2}
\end{align}
$P_k$는 각 단계의 integrated absorbed photons/H다. Full의 incoming은
\begin{equation}
Q_F=\RF Z_0+2bJ.
\label{eq:qfull}
\end{equation}

비광자 source는
\begin{equation}
\BB_k(g,\lambda)
=F_k^{\rm np}(g)+\lambda q_k(g)\nuH-2H_k w\ew,\qquad
\nuH=(1,0,0,-\chi_H)^{\mathsf T}.
\label{eq:bgas}
\end{equation}
실제 gas 단계식은
\begin{equation}
g_1-g_0=h\BB_1(g_1)+\mathsf VP_1,\qquad
g_2-g_1=h\BB_2(g_2)+\mathsf VP_2.
\label{eq:gassteps}
\end{equation}
Full residual의 부호를
\begin{equation}
\HF(g)=g-g_0-2h\BB_2(g)-\mathsf V L_F(g)Q_F,\qquad
\HF(g_F)=0
\label{eq:fullresidual}
\end{equation}
로 고정한다.

비교점 $g_2$의 가상 Full photon은
\begin{equation}
\widehat Z_F=A_F(g_2)Q_F,\qquad
\widehat P_F=L_F(g_2)Q_F
\label{eq:virtualfull}
\end{equation}
다. 실제 Full photon은 $Z_F=A_F(g_F)Q_F$다. 첫 비교점은 residual을 구성하고, 두 번째 비교점은 최종 photon endpoint 차이를 정의한다.

\section{상관된 기체 결함의 정확한 잔차}\label{sec:residual}
\begin{theorem}[같은 COMMON의 residual]\label{thm:residual}
$p=P_1+P_2-\widehat P_F$, $r=\HF(g_2)$라 두면
\begin{equation}
\boxed{r=h[\BB_1(g_1)-\BB_2(g_2)]+\mathsf Vp.}
\label{eq:residualdecomp}
\end{equation}
\end{theorem}
\begin{proof}
식 \eqref{eq:gassteps}를 더하면
$g_2-g_0=h\BB_1(g_1)+h\BB_2(g_2)+\mathsf V(P_1+P_2)$다.
같은 $g_0,Q_F,\lambda,b$를 쓰는 \eqref{eq:fullresidual}에 대입하면 결론을 얻는다. COMMON의 동일성이 항등식 단계에서 반영된다.
\end{proof}

중간 gas 이동도 같은 끝 stage 함수로 결박한다. 다음을 정의한다.
\begin{align}
\delta g_{21}&=g_2-g_1,&
\delta\BB_t(g)&=\BB_1(g)-\BB_2(g),\\
\overline J_{12}
&=\int_0^1D_g\BB_2(g_1+s\delta g_{21})\,\dd s.
\end{align}
미적분학의 기본정리에서
\begin{equation}
\BB_1(g_1)-\BB_2(g_2)
=\delta\BB_t(g_1)-\overline J_{12}\delta g_{21}
\end{equation}
이므로 비광자 residual 부분은
\begin{equation}
\boxed{
r_g=h\delta\BB_t(g_1)
-h\overline J_{12}\bigl[h\BB_2(g_2)+\mathsf VP_2\bigr].
}
\label{eq:gasdecomp}
\end{equation}
첫 항은 source 시각의 계수 변화, 둘째 항은 실제 기체 이동이다. 제\ref{sec:bridge}절에서 첫 시각과 끝 시각의 source를 정확히 연결한다.

이 구성은 독립적인 source interval 둘을 빼는 대신 동일 함수의 유한 차이를 평균 도함수에 결박한다. Interval의 입력 의존성과 이러한 표현 선택의 일반적 배경은 Rump의 제6절, slope 표현은 제11.5절, 비선형 포함은 제13절을 참고할 수 있다.\cite{rump} 아래의 chronology와 source 연결은 이 프로젝트의 단계식에서 직접 유도한다.

\section{광자 absorption 결함의 일곱 작용}\label{sec:photon}
기준점을 $g_2$에 두고
\begin{align}
A&=A_2(g_2),&L&=I-A,&
A_F&=A_F(g_2),&L_F&=I-A_F,\\
\tau&=h\diag\kappa_2(g_2),&
\Delta A_1&=A_1(g_1)-A,&
\Delta\R&=\R^2-\RF
\end{align}
로 쓴다. 대각 연산자
\begin{equation}
C_1=\tau A^2A_F,\qquad C_2=\tau^2A^2A_F
\label{eq:c12}
\end{equation}
는 제\ref{sec:bound}절의 gas 전처리 행렬 $C$와 구별한다. Scalar $t=h\kappa\ge0$에서
\begin{align}
C_1(t)&=\frac{t}{(1+t)^2(1+2t)},&
C_2(t)&=\frac{t^2}{(1+t)^2(1+2t)},\label{eq:c12scalar}\\
A^2-A_F&=-C_2,&A^2+A-2A_F&=C_1.
\label{eq:c12identities}
\end{align}

\begin{theorem}[같은 입력의 integrated absorption 분해]\label{thm:photon}
$[\R,A]=\R A-A\R$라 두면
\begin{equation}
\boxed{\begin{aligned}
p={}&-(I-L\R)\Delta A_1Q_1+L[\R,A]Q_1+L(I-\R)\R Z_0\\
&+C_2\R^2Z_0+L_F\Delta\R Z_0\\
&+bLA(\R-I)J-bC_1J.
\end{aligned}}
\label{eq:seven}
\end{equation}
\end{theorem}
\begin{proof}
$L_1=L-\Delta A_1$을 실제 단계식에 넣으면
\[
p=-(I-L\R)\Delta A_1Q_1+LQ_1+L\R A Q_1+bLJ-L_FQ_F.
\]
$\R A=A\R+[\R,A]$, $Q_1=\R Z_0+bJ$를 대입한다. COMMON 부분에서
\[
L\R Z_0+LA\R^2Z_0-L_F\RF Z_0
=L(I-\R)\R Z_0+C_2\R^2Z_0+L_F\Delta\R Z_0
\]
이고 birth 부분에서
\[
b(2L+LA\R-2L_F)J=bLA(\R-I)J-bC_1J
\]
이다. 두 식은 $L+LA=I-A^2$와 \eqref{eq:c12identities}에서 나온다.
\end{proof}

\begin{table}[H]
\centering\small
\begin{tabularx}{\textwidth}{@{}p{57mm}X@{}}
\toprule
흡수량의 항 & 실제로 분리한 작용\\
\midrule
$-(I-L\R)\Delta A_1Q_1$ & 첫 stage와 끝점 기준 opacity의 변화\\
$L[\R,A]Q_1$ & 중간 흡수와 다음 수송의 순서\\
$L(I-\R)\R Z_0$ & 서로 다른 시점의 흡수를 해당 원 node에서 평가하는 보정\\
$C_2\R^2Z_0$ & 기존 광자의 두 half BE 흡수와 Full BE 흡수\\
$L_F\Delta\R Z_0$ & 같은 COMMON spectrum의 수송 분할\\
$bLA(\R-I)J$ & 먼저 생성된 광자의 추가 redshift\\
$-bC_1J$ & 총량이 같은 두 birth의 흡수 시점\\
\bottomrule
\end{tabularx}
\caption{정리 \ref{thm:photon}의 일곱 sector. 모든 항은 같은 $Z_0,\lambda,b$를 공유한다.}
\label{tab:sectors}
\end{table}

Opacity 차이도 작은 항을 먼저 묶는다.
\begin{align}
\Delta A_{1,j}
&=h[\kappa_{2,j}(g_2)-\kappa_{1,j}(g_1)]A_{1,j}A_j,
\label{eq:dA}\\
\kappa_{2,j}-\kappa_{1,j}
&=c\sigma_j(n_2-n_1)(1-x_2)+a_{1,j}(x_1-x_2).
\label{eq:dkappa}
\end{align}
식 \eqref{eq:dkappa}의 첫 항은 원 captured density 변화이고 둘째 항은 중간 기체 이동이다. 동일 $c,\sigma_j$를 보존하면 두 큰 opacity를 따로 포함하여 빼는 입력 반복을 줄일 수 있다.

\section{원 spectrum에서 얻는 sector의 부호}\label{sec:signs}
\subsection{중간 흡수와 remap 순서}
원 branch에서
\begin{equation}
\R e_j=(1-\alpha_j)e_j+\alpha_je_{j-1},\qquad
\alpha_j=(1-\rho)\frac{E_j}{E_j-E_{j-1}},\quad j\ge1.
\label{eq:alpha}
\end{equation}
따라서 commutator의 retained 비영 성분은
\begin{equation}
[\R,A]_{j-1,j}
=\alpha_j(A_j-A_{j-1})
=\frac{h\alpha_j(\kappa_{2,j-1}-\kappa_{2,j})}
{(1+h\kappa_{2,j})(1+h\kappa_{2,j-1})}.
\label{eq:comm}
\end{equation}
원 group 0은 HI inactive이므로 guard로 나가는 성분의 survival도 1이며, 그 경계 commutator는 0이다.

\begin{proposition}[흡수--수송 순서의 부호]\label{prop:comm}
원 cross section, $Q_1\ge0$, $0\le x_2\le1$에서
\begin{equation}
\boxed{L[\R,A]Q_1\ge0}
\label{eq:commsign}
\end{equation}
가 성분별로 성립한다.
\end{proposition}
\begin{proof}
Captured cross section은 $\sigma_{16}>\sigma_{17}>\cdots>\sigma_{32}>0$이다. Active output으로 내려가는 edge에서는 $\kappa_{j-1}\ge\kappa_j$, 따라서 $A_j-A_{j-1}\ge0$이다. Threshold edge $16\to15$, 다른 inactive 및 guard output에서는 $L_i=0$이다. 식 \eqref{eq:comm}의 나머지 인수는 비음수다.
\end{proof}
$\epsilon_j>0$인 active output에서 $\nu_j$로 사영하므로 이 sector의 직접 $x,w$ forcing은 비음수다.

\subsection{stock과 birth가 경쟁하는 흡수 시점}
두 대각 chronology 항은
\begin{equation}
\boxed{C_2\R^2Z_0-bC_1J=C_1[\tau\R^2Z_0-bJ].}
\label{eq:balance}
\end{equation}
Birth node 24에서
\[
\tau_{24}(\R^2Z_0)_{24}\gtreqless bM
\]
이 이 sector의 흡수 증가·영·감소를 결정한다. 다른 active group에서는 birth 성분이 0이므로 이 항은 비음수다. 특히 $\R=I,A_1=A$에서는 전체 absorbed-photon 차이가
\begin{equation}
p_j=
\frac{\tau_j}{(1+\tau_j)^2(1+2\tau_j)}
(\tau_jN_{0,j}-bJ_j)
\label{eq:scalarbalance}
\end{equation}
로 줄어든다. 같은-opacity 한 group에서 stock의 두 번 흡수 효과와 새 birth의 시점 효과를 직접 비교하는 식이다.

\subsection{먼저 생성된 광자의 redshift}
$\beta_j=\tau_j/(1+\tau_j)^2$, $\eta=hcn_2(1-x_2)$로 두면
\begin{equation}
p_{bR}=bLA(\R-I)J
=bM\alpha_{24}(\beta_{23}e_{23}-\beta_{24}e_{24}).
\label{eq:birthred}
\end{equation}
직접 공통 분모를 취하면
\begin{equation}
\beta_{23}-\beta_{24}
=\frac{\eta(\sigma_{23}-\sigma_{24})(1-\eta^2\sigma_{23}\sigma_{24})}
{(1+\eta\sigma_{23})^2(1+\eta\sigma_{24})^2}.
\label{eq:betadiff}
\end{equation}

\begin{proposition}[birth-redshift의 이온화와 열 방향]\label{prop:birthsign}
원 captured leaves와 branch에서 $b>0,x_2<1$이면
\begin{equation}
\boxed{\boldsymbol1^{\mathsf T}p_{bR}>0,\qquad
\epsilon^{\mathsf T}p_{bR}<0.}
\label{eq:birthsign}
\end{equation}
\end{proposition}
\begin{proof}
전체 $0\le x_2\le1$에서
\[
\eta^2\sigma_{23}\sigma_{24}
\le(hcn_2)^2\sigma_{23}\sigma_{24}
\simeq0.00013495931272495<1.
\]
이 상한과 $\sigma_{23}>\sigma_{24}$는 원 bits의 exact 유리수 비교로 확인되었다. 그러므로 \eqref{eq:betadiff}가 양수다. 또한
$\epsilon_{23}\sigma_{23}<\epsilon_{24}\sigma_{24}$이고
$(1+\eta\sigma_{23})^2>(1+\eta\sigma_{24})^2$이므로
$\epsilon_{23}\beta_{23}<\epsilon_{24}\beta_{24}$다.
$bM\alpha_{24}>0$를 곱하면 두 부호를 얻는다.
\end{proof}

따라서 이 sector는 이온화 absorption을 늘리면서 thermal deposition을 줄인다. 현재 기체의 $w,\Pi>0$에 대한 온도 covector
\begin{equation}
T_g\,\delta g=C_T\left(\frac{\delta w}{\Pi}
-\frac{w\,\delta e}{\Pi^2}\right)
\end{equation}
로 $\mathsf Vp_{bR}$를 사영한 값도 음수다. 전체 endpoint의 온도 기여는 제\ref{sec:observables}절에서 같은 inverse를 통과시켜 정의한다.\cite{photon_theory,checks}

\section{원 비균일 격자의 remap 분할과 경계}\label{sec:remap}
PHYS14에서 $\Delta_j=E_j-E_{j-1}$를 사용하여 다음 exact identity를 얻었다.
\begin{equation}
\boxed{
\Delta\R e_j
=\alpha_j\alpha_{j-1}
\left[
\frac{\Delta_{j-1}}{\Delta_j}e_j
-\left(1+\frac{\Delta_{j-1}}{\Delta_j}\right)e_{j-1}
+e_{j-2}
\right],\quad j\ge2.
}
\label{eq:remapdef}
\end{equation}
이는 $\rho E_j,\rho^2E_j\in[E_{j-1},E_j]$ 및
$\rho E_{j-1}\in[E_{j-2},E_{j-1}]$인 원 branch에 적용된다.
현재의 captured 격자와 기하학에서 이 branch 조건은 PHYS14에 포함되어 있다.

경계 column은 guard를 포함하여
\begin{align}
\Delta\R e_1={}&\alpha_1(\alpha_1+\rho-1)e_1
-\alpha_1(\alpha_1+\rho)e_0\nonumber\\
&+\alpha_1e_{G_N}+\rho E_0\alpha_1e_{G_U}.
\label{eq:boundarydef}
\end{align}
$j=0$과 기존 두 guard column의 차이는 0이다. 식
\eqref{eq:remapdef}--\eqref{eq:boundarydef}의 계수합 및 energy 가중합은 정확히 0이므로
\begin{equation}
\mathfrak n\Delta\R=0,\qquad \mathfrak u\Delta\R=0.
\label{eq:zero_moments}
\end{equation}

따라서 remap 두 경로가 보존하는 number와 redshifted energy가 같아도 energy node별 분포가 달라진다. 제\ref{sec:photon}절에서는 이 spectrum 차이를 흡수 연산자 $L_F$로 전달한다. 원 captured node의 작은 비균일성과 아래 경계 수지는 그대로 유지한다. 새로운 interval 계산에서도 $\Delta\R$를 큰 두 결과의 차이로 만드는 대신 위 연결된 계수에서 구성한다.\cite{phys14,photon_theory}


\section{원 source의 밀도와 온도 정규화 연결}\label{sec:bridge}
Full과 둘째 half는 같은 끝 시각의 source를 쓰지만 첫 half의 저장 density는 다르다. 이 절에서는
\begin{equation}
f_k=\fh_k=\frac{\ex(n_{{\rm He},k}^{\rm stored})}{\ex(n_k)},\qquad
s=y_1+2y_2,\qquad e_k=x+f_ks
\end{equation}
로 쓴다. 전자 수와 전체 자유입자 수, 온도는
\begin{equation}
\Pi_k=1+f_k+e_k,\qquad
T_k=C_T\frac{w}{\Pi_k},\qquad
C_T=\frac{2\ex(\epsilon_{\rm eV})}{3\ex(k_B)}.
\label{eq:thermo}
\end{equation}
$f_k$는 이 절의 source density ratio이고, 원 stage에 입력된 scalar $f_{\rm He,stage}$와의 대응은 저장 곱 $n_kf_{\rm He,stage}$를 통해 정해진다.

\subsection{FT03 반응 함수의 완전한 인수 분해}
종 표기를 $a=H,I,II$로 두며 $I,II$는 각각 HeI, HeII 이온화 채널이다.
$\beta_a(T)$는 collisional-ionization coefficient,
$\alpha_a(T)$는 radiative-recombination coefficient,
\[
s_a(T)=\frac{\dd\log\alpha_a}{\dd\log T}
\]
이다. 모든 함수는 원 FT03의 고정 fit이며 계수의 단위는 $\mathrm{cm^3\,s^{-1}}$다. HeI dielectronic recombination은
\begin{equation}
d_{\rm DR}(T)=d_1(T)+d_2(T),\qquad
d_r(T)=c_rA_{\rm DR}T^{-3/2}\exp(-B_r/T),\quad(c_1,c_2)=(1,0.3).
\label{eq:dr}
\end{equation}
$A_{\rm DR},B_r$는 원 source에 저장된 세 고정 leaves다.

분율의 반응 인수는 다음과 같이 모두 주어진다.
\begin{align}
k_x&=(1-x)\beta_H-x\alpha_H,\label{eq:kx}\\
k_y&=(1-y_1-y_2)\beta_I-y_1(\alpha_I+d_{\rm DR})
-y_1\beta_{II}+y_2\alpha_{II},\label{eq:ky}\\
k_z&=y_1\beta_{II}-y_2\alpha_{II},\qquad
k_{\rm frac}=(k_x,k_y,k_z)^{\mathsf T}.
\label{eq:kz}
\end{align}
Photo-off source의 collisional·RR·DR 항은 반응 전 population과 electron density의 곱이다. H당 정상화 및 He 분율의 정상화를 적용하면
\begin{equation}
\boxed{F^{\rm np}_{k,{\rm frac}}(g)
=n_ke_k\,k_{\rm frac}(x,y_1,y_2,T_k).}
\label{eq:npfrac}
\end{equation}

RR의 kinetic energy coefficient를
\begin{equation}
\mathcal K_a(T)=
\frac{k_BT\,\alpha_a(T)\,[3/2+s_a(T)]}{\epsilon_{\rm eV}}
\label{eq:kinetic}
\end{equation}
로 둔다. $\mathcal K_a$의 단위는 $\mathrm{eV\,cm^3\,s^{-1}}$다.
FT03 controlled model의 threshold를 $\chi_a^{\rm FT}$로 쓰면
\begin{equation}
k_{wH}=-\chi_H^{\rm FT}(1-x)\beta_H-x\mathcal K_H,
\label{eq:kwH}
\end{equation}
\begin{align}
k_{w{\rm He}}={}&-\chi_I^{\rm FT}(1-y_1-y_2)\beta_I
-\chi_{II}^{\rm FT}y_1\beta_{II}
-y_1\mathcal K_I-y_2\mathcal K_{II}\nonumber\\
&-y_1\sum_{r=1}^2d_r(T)\frac{k_BB_r}{\epsilon_{\rm eV}}.
\label{eq:kwHe}
\end{align}
따라서 열 source는 정확히
\begin{equation}
\boxed{F^{\rm np}_{k,w}(g)
=n_ke_k\bigl[k_{wH}(x,T_k)+f_k k_{w{\rm He}}(y_1,y_2,T_k)\bigr].}
\label{eq:npthermal}
\end{equation}
이 $F^{\rm np}$는 HH·photo·팽창을 넣기 직전에 export한 네 성분이다. 팽창은 \eqref{eq:bgas}에서 따로 더한다.

이번 controlled model의 실제 정의 사슬을 따라가면
\begin{equation}
(\chi_H^{\rm FT},\chi_I^{\rm FT},\chi_{II}^{\rm FT})
=\ex(13.598434599702,\ 24.587389011,\ 54.41776)
\label{eq:thresholds}
\end{equation}
이며 native photon/HH의 세 $\chi_a$ leaves와 bits가 각각 같다.
따라서 이 고정 source에서는 $\chi_a^{\rm FT}=\chi_a$로 동일시할 수 있다.\cite{bindings,threshold_binding}

HH의 별도 event rate는 원 source에서
\begin{equation}
q_k(g)=n_k(1-x)^2 k_{\rm HH}(T_k),\qquad
k_{\rm HH}(T)=\ex(1.2\times10^{-17})T^{\ex(1.2)}
\exp[-\ex(157800)/T]
\label{eq:hh}
\end{equation}
다. 온도 정의는 \eqref{eq:thermo}와 같다.

\subsection{온도를 맞추는 사상}
\begin{theorem}[source 시각의 정확한 bridge]\label{thm:bridge}
같은 분율에서
\begin{equation}
\Psi g=(x,y_1,y_2,w\Pi_2/\Pi_1),\quad
r_n=\frac{n_1}{n_2},\quad
r_e=r_n\frac{e_1}{e_2},\quad \Delta f=f_1-f_2
\label{eq:psi}
\end{equation}
로 두면 $T_2(\Psi g)=T_1(g)$이고
\begin{align}
F^{\rm np}_1(g)
&=r_e F^{\rm np}_2(\Psi g)
+n_1e_1\Delta f\,k_{w{\rm He}}(y_1,y_2,T_1)\ew,\label{eq:bridgeF}\\
q_1(g)&=r_nq_2(\Psi g).
\label{eq:bridgeq}
\end{align}
\end{theorem}
\begin{proof}
$\Psi$는 분율과 $\Pi_2$를 고정하면서 $w$를 $w\Pi_2/\Pi_1$로 바꾸므로
$C_T\Psi_w/\Pi_2=C_Tw/\Pi_1$이다.
따라서 \eqref{eq:npfrac}의 모든 $k_{\rm frac}$ 함수가 같은 온도에서 평가된다.
열 항은
\[
r_e F^{\rm np}_{2,w}(\Psi g)
=n_1e_1[k_{wH}(T_1)+f_2k_{w{\rm He}}(T_1)]
\]
이므로 $f_1-f_2$에 비례한 항을 더하면 \eqref{eq:bridgeF}다.
식 \eqref{eq:hh}에서는 온도를 맞춘 뒤 남는 비율이 $n_1/n_2$이므로
\eqref{eq:bridgeq}를 얻는다.
\end{proof}

\subsection{작은 차이를 먼저 구성한다}
비율과 thermal-coordinate 이동을
\begin{equation}
r_e-1=(r_n-1)+r_n\Delta f\,\frac{s}{e_2},\qquad
\delta w_\Psi=\Psi_w(g)-w=-w\Delta f\,\frac{1+s}{\Pi_1}
\label{eq:small_bridge}
\end{equation}
로 구성한다. $g$에서 $\Psi g$로 가는 $w$-선분의 평균 derivative를
$\overline{F^{\rm np}_{2,w}}$, $\overline{q_{2,w}}$로 쓰면
\begin{align}
F^{\rm np}_1(g)-F^{\rm np}_2(g)
={}&(r_e-1)F^{\rm np}_2(g)
+r_e\overline{F^{\rm np}_{2,w}}\,\delta w_\Psi\nonumber\\
&+n_1e_1\Delta f\,k_{w{\rm He}}(T_1)\ew,\label{eq:timeF}\\
q_1(g)-q_2(g)
={}&(r_n-1)q_2(g)
+r_n\overline{q_{2,w}}\,\delta w_\Psi .
\label{eq:timeq}
\end{align}
원 $H_1=H_2$여서 같은 $g$의 팽창 항은 $\delta\BB_t$에서 소거된다.
식 \eqref{eq:timeF}--\eqref{eq:timeq}와 $\lambda\nuH$를 결합하면
\eqref{eq:gasdecomp}의 source 시각 항이 된다.

원 exact leaves에서
\begin{equation}
\Delta f\in[-6.94524945677013820,\,-6.94524945677013819]\times10^{-18}.
\label{eq:dfvalue}
\end{equation}
Full·first-half·실제 second-half root image의 공동 convex hull에서는
\begin{equation}
\delta w_\Psi\in[1.09800931589085210,\,
1.11069895932571596]\times10^{-16}\ {\rm eV/H}.
\label{eq:dwpsivalue}
\end{equation}
이 hull과 $\Psi$의 상, 그 사이의 선분이 원 $X$ 내부임을 exact endpoint로 확인했다.
따라서 기존 whole-$X$ gradient·Hessian 포함을 이 연결에 적용할 수 있다.
이와 같이 실제 사용할 내부 영역을 지정하면 source를 다시 호출하지 않고 원 derivative 근거를 재사용한다.\cite{bindings}

\section{기존 Full 포함으로 닫히는 endpoint 결함}\label{sec:bound}
\subsection{평균 Jacobian과 photo 한 열}
\begin{theorem}[유한 endpoint 전달]\label{thm:transfer}
$d=g_2-g_F$와
\begin{equation}
\overline{\mathcal A}
=\int_0^1D_g\HF(g_F+s d)\,\dd s
\label{eq:Abar}
\end{equation}
에 대해
\begin{equation}
\boxed{\overline{\mathcal A}d=r.}
\label{eq:mean_defect}
\end{equation}
\end{theorem}
\begin{proof}
원 convex $X$가 $g_F,g_2$와 연결 선분을 포함한다.
미적분학의 기본정리를 적용하면
$\HF(g_2)-\HF(g_F)=\overline{\mathcal A}(g_2-g_F)$다.
왼쪽은 정의에 의해 $r$이다.
\end{proof}

Photo의 current-gas derivative는 $x$ 한 열만 가진다.
\begin{align}
D_g\HF(g)
&=I-2hD_g\BB_2(g)+2h\,v_F(g)e_x^{\mathsf T},\label{eq:photojac}\\
v_F(g)
&=\sum_j
\frac{a_{2,j}Q_{F,j}}{[1+2h\kappa_{2,j}(g)]^2}\nu_j.
\label{eq:vf}
\end{align}
여기서 $Q_F$는 current-gas unknown에 독립인 공통 입력이다. 원 Full source 포함은 같은 $Q_F$와 같은 계수 계열에서 \eqref{eq:photojac}를 포함한다.

\subsection{고정 전처리 행렬과 componentwise bound}
원 NCP07 Full의 고정 binary64 행렬 $C$를 exact dyadic 행렬로 사용한다.
$C$는 centre Jacobian에서 계산되어 저장된 approximate inverse다. 원 증명의 scale과 포함은
\begin{gather}
S=\diag\bigl(\ex(0.015),\ex(0.01),\ex(0.01),\ex(0.5)\bigr),
\label{eq:S}\\
|I-CD_g\HF(g)|\le B\quad(g\in X),\qquad
\widetilde B=S^{-1}BS,\qquad
\normi{\widetilde B}\le q<1,\label{eq:Bbound}\\
q=\ex(0.00011736830563768177).
\label{eq:q}
\end{gather}
$C,B,q,S$는 같은 원 \texttt{UNIFORM\_ARITHMETIC.json}에서 결박한다.
표시가 매우 가까운 다른 diagnostic export와 섞지 않는다.\cite{preconditioner}

\begin{theorem}[재사용된 Full bound로부터의 결함 포함]\label{thm:majorant}
원 포함 \eqref{eq:Bbound} 아래 $\overline{\mathcal A}$는 가역이며
\begin{equation}
\boxed{
|S^{-1}d|\le(I-\widetilde B)^{-1}|S^{-1}Cr|,
\qquad
\normi{S^{-1}d}\le\frac{\normi{S^{-1}Cr}}{1-q}.
}
\label{eq:majorant}
\end{equation}
\end{theorem}
\begin{proof}
$\overline E=I-C\overline{\mathcal A}$라 두면 평균 적분에서도
$|\overline E|\le B$다. 따라서
$\normi{S^{-1}\overline ES}\le q<1$이고 $I-\overline E=C\overline{\mathcal A}$가 가역이다.
유한 정사각 행렬의 곱이 가역이므로 $C,\overline{\mathcal A}$도 가역이다.
\[
d=Cr+\overline E d,\qquad
S^{-1}d=\sum_{m=0}^{\infty}(S^{-1}\overline ES)^mS^{-1}Cr.
\]
각 항의 절대값을 비음수 $\widetilde B^m$로 지배하여 첫 식을 얻고,
노름의 기하급수를 합하면 둘째 식이다.
\end{proof}

원 $q$의 scalar amplification은
\begin{equation}
\frac1{1-q}<1.000118.
\label{eq:amplification}
\end{equation}
이 상한은 \emph{전처리된 잔차 $Cr$}의 원 반경 척도 노름에 적용된다.
실제 residual의 크기는 \eqref{eq:residualdecomp},
\eqref{eq:seven}, \eqref{eq:timeF}--\eqref{eq:timeq}를 같은 입력에서 평가해 정한다.
평균 Jacobian을 끝점 Jacobian으로 대체하는 경우의 차이는 별도의
비선형 판별 예제에서 확인했다.

포함을 계산할 때는 먼저 \eqref{eq:majorant}에서 유효한 초기 $[d]$를 만들고,
원 $[E]=I-C[D\HF](X)$를 사용하여
\begin{equation}
[d]\leftarrow[d]\cap\bigl(C[r]+[E][d]\bigr)
\label{eq:refine}
\end{equation}
를 적용할 수 있다. 이미 있는 $[g_2]-[g_F]$도 독립적인 valid cap으로 교차한다.
더 좁은 $[g_2]-[d]$와 source Hessian을 활용할지는 실제 남은 폭에 따라 정한다.
평균 Jacobian을 통한 포함의 일반적 원리는 비선형 검증 이론을 따르며,
이곳에서는 고정된 원 $C,B,q,S$와 실제 두 단계 residual을 결박했다.\cite{rump}

\section{유한 관측량과 같은 inverse의 sector 기여}\label{sec:observables}
두 endpoint는 같은 $f_2,C_T,n_2$를 사용한다. 전자 covector와 열 좌표는
\begin{equation}
\ell_e=(1,f_2,2f_2,0),\qquad
d_e=\ell_e d,\qquad d_w=e_w^{\mathsf T}d,\qquad
\Delta n_e=n_2d_e.
\label{eq:electron_defect}
\end{equation}
$\Pi_F=\Pi_2-d_e>0$이고 $T_2=C_Tw_2/\Pi_2$이다.

\begin{theorem}[유한 온도 차이와 공통 adjoint]\label{thm:temperature}
정확한 유한 온도 차이는
\begin{equation}
\boxed{\Delta T=T_2-T_F
=\frac{C_Td_w-T_2d_e}{\Pi_2-d_e}.}
\label{eq:finiteT}
\end{equation}
또한
\begin{equation}
\ell_T=\frac{C_Te_w^{\mathsf T}-T_2\ell_e}{\Pi_F},\qquad
\overline{\mathcal A}^{\mathsf T}z_T=\ell_T^{\mathsf T}
\label{eq:adjoint}
\end{equation}
이면
\begin{equation}
\boxed{\Delta T=z_T^{\mathsf T}r,\qquad
r=\sum_s r_s\ \Longrightarrow\
\Delta T=\sum_s z_T^{\mathsf T}r_s.}
\label{eq:attribution}
\end{equation}
\end{theorem}
\begin{proof}
$T_F=C_T(w_2-d_w)/(\Pi_2-d_e)$를 $T_2$에서 빼고 공통 분모를 취하면
\eqref{eq:finiteT}다. 이는 $\Delta T=\ell_Td$와 같다.
식 \eqref{eq:mean_defect}를 이용하여
$\ell_Td=z_T^{\mathsf T}\overline{\mathcal A}d=z_T^{\mathsf T}r$를 얻는다.
같은 $z_T$에 대한 선형성이 마지막 합을 준다.
\end{proof}

$r_s$는 source 시각·중간 gas 이동 및 광자 일곱 sector다.
모든 항에 같은 원 endpoint와 같은 평균 inverse가 적용된다.
따라서 각 기여는 하나의 유한 $\Delta T$를 정확히 분할한다.
전자 수, 열에너지, 전자 밀도도 대응 covector를 동일한 방식으로 사용한다.

식 \eqref{eq:finiteT}는 큰 두 온도 포함을 각각 뺄 필요 없이
같은 $d$에서 $C_Td_w-T_2d_e$를 구성한다.
제\ref{sec:signs}절의 직접 forcing 부호를 endpoint 기여로 변환하는 대상은
바로 \eqref{eq:attribution}다.
이온화로 늘어난 자유입자와 열에너지 변화가 온도에서 결합하는 방식도
유한 분모 $\Pi_F$를 포함해 유지된다.

\section{guard ledger와 H 화학·열에너지 잔차}\label{sec:ledger}
잔차 기준 photon 차이와 그 moments를
\begin{equation}
\Delta\widehat Z=Z_2-\widehat Z_F,\qquad
\Delta\mathcal N=\mathfrak n\Delta\widehat Z,\qquad
\Delta\mathcal U=\mathfrak u\Delta\widehat Z
\label{eq:photon_moments}
\end{equation}
로 쓴다. 두 경로의 총 birth가 같고 $\mathfrak n\R=\mathfrak n$이므로
\begin{equation}
\boldsymbol1^{\mathsf T}p=-\Delta\mathcal N.
\label{eq:number_ledger}
\end{equation}

\begin{theorem}[흡수 chronology와 energy residual]\label{thm:energy}
원 guard 수지에서
\begin{equation}
\boxed{E\cdot p
=(1-\rho)(E\cdot P_1-bME_{24})-\Delta\mathcal U.}
\label{eq:energy_ledger}
\end{equation}
\end{theorem}
\begin{proof}
첫 half의 absorption은 두 번째 redshift 전에 발생한다. 식
\eqref{eq:remapmom}을 순서대로 적용하면
\begin{align*}
\mathfrak uZ_2
&=\rho^2\mathfrak uZ_0+b(\rho+1)ME_{24}
-\rho E\cdot P_1-E\cdot P_2,\\
\mathfrak u\widehat Z_F
&=\rho^2\mathfrak uZ_0+2bME_{24}-E\cdot\widehat P_F.
\end{align*}
두 식을 빼고 $p=P_1+P_2-\widehat P_F$를 대입한다.
\end{proof}

원 active 17개 group의 captured heat는
\begin{equation}
\epsilon_j=\ex(E_j)-\ex(\chi_H)
\label{eq:heat_identity}
\end{equation}
와 정확히 같다. 실제 bits에서 확인했고,
$\chi_H\le E_j<2\chi_H$인 normal binary subtraction의 Sterbenz 조건과도 일치한다.
원 subtraction capture를 유지한 상태에서 이 identity를 사용할 수 있다.
Inactive group의 흡수는 0이므로
\begin{equation}
\epsilon\cdot p
=(1-\rho)(E\cdot P_1-bME_{24})
-\Delta\mathcal U+\chi_H\Delta\mathcal N.
\label{eq:thermal_ledger}
\end{equation}

H 화학·열에너지 covector
\begin{equation}
\ell_H=(\chi_H,0,0,1)
\end{equation}
는 $\ell_H\nu_j=E_j$, $\ell_H\nuH=0$를 만족한다.
따라서
\begin{equation}
\boxed{\begin{aligned}
\ell_Hr={}&h\ell_H[\BB_1(g_1)-\BB_2(g_2)]\\
&+(1-\rho)(E\cdot P_1-bME_{24})-\Delta\mathcal U.
\end{aligned}}
\label{eq:Henergy}
\end{equation}
이 사영의 물리량은 $\chi_Hx+w$다.
HH의 직접 $q\nuH$는 소거되고, HH가 실제 기체·opacity·흡수·nonphoto 평가점에 전달한 효과가 나머지 항에 남는다.
Helium의 nonphoto 열 기여도 원 벡터 안에 유지된다.

$(1-\rho)E\cdot P_1$은 첫 absorption이 이미 소모한 광자에 대한 chronology 보정이고,
$-(1-\rho)bME_{24}$는 먼저 주입된 birth가 겪는 redshift를 기록한다.
따라서 이 ledger는 spectrum 모양과 별개로 시간 순서를 점검하는 정확한 관계다.

\section{실제 Full photon endpoint의 회수}\label{sec:actual_photon}
같은 $g_2$를 기준으로 한 outgoing 차이는
\begin{equation}
\begin{aligned}
\Delta\widehat Z={}&A\R\Delta A_1Q_1+A[\R,A]Q_1-C_2\R^2Z_0
+A_F\Delta\R Z_0\\
&+bA^2(\R-I)J+bC_1J.
\end{aligned}
\label{eq:outgoing}
\end{equation}
이는 \eqref{eq:stage1}--\eqref{eq:virtualfull}에
$\R A=A\R+[\R,A]$와 \eqref{eq:c12identities}를 적용한 정확한 항등식이다.

실제 Full과의 차이 $\Delta Z=Z_2-Z_F$는 같은 $d_x=x_2-x_F$를 사용하여 회수한다.
\begin{equation}
\boxed{
\Delta Z_j=\Delta\widehat Z_j+
\frac{2h\,a_{2,j}Q_{F,j}\,d_x}
{D_{F,j}(g_2)D_{F,j}(g_F)},\qquad 0\le j\le32,
}
\label{eq:fullphotonsecant}
\end{equation}
\begin{equation}
D_{F,j}(g)=1+2h\kappa_{2,j}(g).
\end{equation}
Guard의 추가항은 0이다.
이 식의 부호는 survival의 다음 resolvent 차이에서 정해진다.
\[
\bigl(A_F(g_2)-A_F(g_F)\bigr)_{jj}
=\frac{2h\,a_{2,j}(x_2-x_F)}{D_{F,j}(g_2)D_{F,j}(g_F)}.
\]

일반 비교점 $g_*$를 기준으로 photon 분해를 만들 수도 있다.
그 경우 실제 둘째 half로 옮기는 추가항은 retained 성분마다
\begin{equation}
\bigl([A_2(g_2)-A_2(g_*)]Q_2\bigr)_j
=\frac{h\,a_{2,j}Q_{2,j}(x_2-x_*)}
{D_{2,j}(g_2)D_{2,j}(g_*)},\quad 0\le j\le32,
\label{eq:halfphotonsecant}
\end{equation}
이며 guard는 0이다.
식 \eqref{eq:fullphotonsecant}는 Full 분모·$Q_F$·$2h$를,
\eqref{eq:halfphotonsecant}는 half 분모·$Q_2$·$h$를 사용한다.
두 비교점의 차이를 명시하면 residual에서 실제 photon endpoint까지 같은 gas 결함을 전달할 수 있다.

\section{판별 검산, 연구 판정과 다음 계산}\label{sec:verification}
\begin{table}[H]
\centering\small
\begin{tabularx}{\textwidth}{@{}p{31mm}X p{35mm}@{}}
\toprule
검산 & 실제 판별 대상 & 결과\\
\midrule
광자 exact Fraction 및 네 슬롯 & outgoing·흡수 분해, opacity secant, remap/guard 수지, transparent·stock·birth 극한 & 39개 관계, 3148개 scalar equality 통과\\
원 leaf 부호 & HI $\sigma$의 순서, active heat identity, birth-redshift의 전체 물리 범위 조건 & exact 유리수 조건 통과\\
비선형 전달 예제 & 비대각 4D quadratic residual, 평균 Jacobian과 끝점 Jacobian의 차이, Neumann bound, finite $T$, 같은 adjoint의 sector 합 & 14개 관계, 16개 scalar equality 통과\\
Wolfram 대수 & $A^2-A_F$, $A^2+A-2A_F$, finite $T$, $\beta_{23}-\beta_{24}$의 factorization & 네 차이 $\{0,0,0,0\}$\\
독립 검토자 & 별도 세 제조 대수 계열, 원 source graph·leaf·고정 $C$, 전체 논증과 타입 & 26개 검토 항목 통과\\
\bottomrule
\end{tabularx}
\caption{실제로 수행한 표적 검산. 관계·scalar equality·검토 항목은 각각 다른 집계 단위다.}
\label{tab:checks}
\end{table}
광자 대수 검산의 wall time은 약 0.48초, 별도 sector leaf 검사는 약 0.04초,
비선형 전달 예제는 약 0.0013초, 독립 검토 실행은 약 0.051초였다.
제조 상태는 항등식과 판별 사례를 확인하는 입력으로 기록했다.
이번 이론 단위의 새 physical source callback, root solve, history 실행은 모두 0회다.
PHYS09--14의 원 source·근 포함은 전제를 유지한 채 재사용했다.\cite{checks,review}

표적 비선형 예제에서는 평균 Jacobian이 정확한 $+r$ 식을 만족하는 반면,
끝점 Jacobian으로 치환한 식에 비영 quadratic remainder가 남았다.
Finite temperature 식과 tangent 식도 구별되었다.
이 판별은 제\ref{sec:bound}--\ref{sec:observables}절에서 유지한 평균과 유한 분모의 역할을 확인한다.



이번 독립 판정은 \textbf{\ReviewDecision}이다.
같은 COMMON, linked geometric isotropic source-aligned $C^2$ 계열,
captured 원 상수와 두 half에 공통인 $\lambda,b$ 아래의 결함 구성,
source bridge, 평균 Jacobian 전달, finite readout, ledger와 부호 정리를 채택했다.
최종 검토에서 제기된 photon secant의 행렬·벡터 표기는
\eqref{eq:fullphotonsecant}--\eqref{eq:halfphotonsecant}의 성분식으로 교정하여 확인했다.\cite{review}

\subsection{다음 포함 계산의 연결}
후속 단위는 \textbf{HH\_CORRELATED\_ENDPOINT\_ENCLOSURE}다.
입력은 같은 COMMON의 원 근·carry·source 포함,
정확한 $C,S$, 그리고 이 보고서의 expression 관계다.
이 단위의 \texttt{assemble\_correlated\_defect}와
\texttt{enclose\_endpoint\_defect}는 아래 식을 구현할 후속 함수의 계산계약이다.
먼저 \eqref{eq:residualdecomp}의 $[r]$를 구성하고,
\eqref{eq:majorant}--\eqref{eq:refine}에서 $[d]$를 얻는다.
이후 \eqref{eq:finiteT}, \eqref{eq:electron_defect},
\eqref{eq:attribution}, \eqref{eq:fullphotonsecant}로 관측량과 sector 기여를 포함한다.

\begin{figure}[H]
\centering
\begin{tikzpicture}[
 box/.style={draw=MidnightBlue!65,fill=MidnightBlue!3,rounded corners=2pt,
 text width=40mm,minimum height=11mm,align=center,font=\small},
 >={Stealth[length=2mm]},node distance=10mm and 9mm]
\node[box] (input) {같은 COMMON\\원 근·source·carry};
\node[box,below left=of input] (gas) {source 시각·기체 이동\\식 \eqref{eq:gasdecomp}, \eqref{eq:timeF}};
\node[box,below right=of input] (photon) {광자 일곱 sector\\식 \eqref{eq:seven}};
\node[box,below=26mm of input] (res) {상관된 residual $r$\\식 \eqref{eq:residualdecomp}};
\node[box,below=of res] (defect) {기존 $C,B,q,S$로 $d$ 포함\\식 \eqref{eq:majorant}};
\node[box,below left=of defect] (readout) {$\Delta T,\Delta e,\Delta w,\Delta n_e$\\같은 adjoint의 기여};
\node[box,below right=of defect] (out) {실제 $\Delta Z$\\동일 $d_x$의 secant};
\draw[->] (input)--(gas);
\draw[->] (input)--(photon);
\draw[->] (gas)--(res);
\draw[->] (photon)--(res);
\draw[->] (res)--(defect);
\draw[->] (defect)--(readout);
\draw[->] (defect)--(out);
\end{tikzpicture}
\caption{이번 이론에서 다음 실제 포함 계산으로 이어지는 의존관계. 같은 원 입력과 동일한 결함이 두 종류의 endpoint 출력을 연결한다.}
\label{fig:dag}
\end{figure}

\subsection{계산 중 유지할 정확한 관계}
\begin{enumerate}
\item $\RF=\R(\rho^2)$, $M_F=2M$, 같은 $Z_0,\lambda,b$를 먼저 결박한다.
\item $\Delta\R$, $[\R,A]$, $\Delta A_1$, $r_e-1$, $\delta w_\Psi$를 차이의 factorized 식에서 구성한다.
\item 원 33개 group, 두 guards와 signed $U,V,W$를 유지한다. 구조적 관계는 원 계수와 source identity로 적용한다.
\item 직접 photon 합과 number/energy ledger의 유효 포함을 교차할 수 있다. 정확한 zero moment는 연결된 계수에서 구성한다.
\item 새로 얻은 실제 폭이 source·remap·birth 중 어느 관계에 의해 지배되는지 읽고, 그 관계에 필요한 포함만 다듬는다.
\end{enumerate}
이론적으로 남아 있던 연결은 실제 두 단계식에서 같은 COMMON의 잔차를 만들고,
이를 기존 Full 증명과 유한 관측량으로 전달하는 사슬이었다.
정리 \ref{thm:residual}--\ref{thm:energy}와 photon secant가 그 사슬을 완성한다.
후속 계산은 이 사슬의 크기와 부호를 실제 원 입력에서 평가한다.


\appendix
\section{기호와 원자료의 대응}
\begin{longtable}{@{}p{30mm}p{115mm}@{}}
\toprule
기호 & 정의와 원자료 대응\\
\midrule
\endhead
$g_0,Z_0$ & 같은 COMMON의 gas 및 aggregate photon/guard. Full·first-half binary payload 대응을 binding에서 확인했다.\\
$g_1,g_2,g_F$ & first half, 실제 둘째 half, 같은 종료 시각 Full의 root family. 원 $X$ 안의 포함을 재사용한다.\\
$Q_k,Z_k,P_k$ & absorption 전, 후, integrated absorbed 상태. $P_k=(I-A_k)Q_k$.\\
$\R,\RF$ & 원 33-node hat remap와 두 guard를 포함한 수송. $\rho$와 $\rho^2$를 결박한다.\\
$A,L,A_F,L_F$ & 기준점 $g_2$의 half/Full survival 및 absorption 대각 연산자. Guard에서는 $A=1,L=0$.\\
$\Delta A_1$ & $A_1(g_1)-A_2(g_2)$, 원 density 차이와 $x_1-x_2$로 factorize한다.\\
$C_1,C_2$ & 광자 chronology의 비음수 대각 연산자. Gas 전처리 행렬 $C$와 별개다.\\
$\HF,r,d$ & Full residual, $r=\HF(g_2)$, 실제 gas 차이 $d=g_2-g_F$.\\
$C,B,q,S$ & 원 Full ARITHMETIC export의 고정 전처리 행렬, magnitude bound, 수축 상한, radius scale.\\
$f_k,\Psi$ & $f_k=n_{{\rm He},k}^{\rm stored}/n_k$와 온도 정렬 사상 $\Psi_w=w\Pi_2/\Pi_1$.\\
$\ell_e,\ell_T,\ell_H$ & 전자 수, 유한 온도 차이, H 화학·열에너지의 covector.\\
$\Delta\widehat Z,\Delta Z$ & 각각 $g_2$ 기준 가상 Full과의 photon 차이, 실제 $g_F$ Full과의 차이.\\
\bottomrule
\end{longtable}

원 threshold의 bits는
\[
(\texttt{402b3266051507a6},\ \texttt{4038965f20504658},\
\texttt{404b357928e0c9da})
\]
이며 FT03 controlled threshold와 photon/HH constants가 각각 동일하다.
원 $C$의 16개 entries, $B$의 16개 bounds, $S$ 및 $q$는
\texttt{FIXED\_PRECONDITIONER\_BINDING.json}의 bits/exact 값을 계산 입력으로 사용한다.
본문의 소수 표시는 이 값을 읽기 쉽게 제시한 것이다.

\section{원고와 검산의 재현 범위}
보고서의 주요 근거는 다음 문서와 JSON에 대응한다. 모든 내부 경로는 PHYS15 연구 묶음의 상대 경로다.
\begin{table}[H]
\centering\small
\begin{tabularx}{\textwidth}{@{}p{41mm}X@{}}
\toprule
내용 & 상대 경로\\
\midrule
통합 이론 & \texttt{theory/CORRELATED\_ENDPOINT\_THEORY\_KO.md}\\
광자 분해·부호·ledger & \texttt{theory/PHOTON\_DEFECT\_THEOREM\_KO.md}\\
source·COMMON 결박 & \texttt{theory/ASSUMPTIONS\_AND\_BINDINGS\_KO.md}\\
고정 $C,B,q,S$ & \texttt{theory/FIXED\_PRECONDITIONER\_BINDING.md}\\
광자 exact 검사 & \texttt{checks/photon/PHOTON\_EXACT\_CHECKS.json}\\
sector 조건 & \texttt{checks/photon/SECTOR\_SIGN\_CHECKS.json}\\
비선형·Wolfram 검사 & \texttt{checks/NONLINEAR\_TRANSFER\_CHECKS.json}, \texttt{checks/WOLFRAM\_ALGEBRA.json}\\
독립 최종 판정 & \texttt{review/DECISION.json}, \texttt{review/REVIEW\_KO.md}\\
\bottomrule
\end{tabularx}
\caption{완성된 이론과 실제 검사 근거의 대응.}
\end{table}

이 보고서는 XeLaTeX으로 컴파일한다. 본문, 표, 정리, TikZ 의존관계 그림 및 한국어 fonts를 source 묶음에 포함한다. 최종 독립 판정 문구는 \texttt{review\_status.tex}, 검산 서술은 \texttt{checks\_status.tex}에 분리되어 있다. 판정의 원자료는 별도의 독립 reviewer 기록으로 유지한다.

\begingroup
\small
\setlength{\parskip}{0pt}
\begin{thebibliography}{9}
\setlength{\itemsep}{4pt}
\bibitem{rump}
S. M. Rump, \emph{Verification methods: Rigorous results using floating-point arithmetic},
Acta Numerica \textbf{19} (2010), 287--449.
일반 배경: \S6, \S11.5, \S13.
\url{https://www.tuhh.de/ti3/paper/rump/Ru10.pdf}.

\bibitem{phys14}
WU088\_HH PHYS14, \emph{공동 carry의 두 번째 half 전달과 광자 기억의 온도 응답},
2026-10-12. 원 승인 commit
\texttt{087a16f28d0ee4fc7ae74ae50bfcafe578716e85}.
\href{https://github.com/cosmosapjw-quantum/WU088_HH/pull/49}{GitHub PR49}.
원 source-aligned closure, 실제 carry 및 geometry/root 근거.

\bibitem{root_theory}
WU088\_HH PHYS15, \emph{동일 COMMON의 Full--two-half 결함 이론},
\texttt{theory/CORRELATED\_ENDPOINT\_THEORY\_KO.md}, 2026-10-13 KST.

\bibitem{photon_theory}
WU088\_HH PHYS15 광자 분과, \emph{공통 입력을 보존하는 광자 결함과 기체 잔차},
\texttt{theory/PHOTON\_DEFECT\_THEOREM\_KO.md}.
Secant 성분 표기 교정 및 원문 보존 기록 포함.

\bibitem{bindings}
WU088\_HH PHYS15, \emph{동일 COMMON endpoint 비교의 전제와 정확한 source 연결},
\texttt{theory/ASSUMPTIONS\_AND\_BINDINGS\_KO.md}와 동반
\texttt{BINDINGS.json}. NCP07 원 archive SHA256:
\texttt{4ea6b7dbbc714493ac85f944a65a353a1ee75abf5742c3e7ce65a7df56cc3465}.

\bibitem{preconditioner}
WU088\_HH PHYS15, \emph{기존 Full preconditioner와 수축 상한의 결박},
\texttt{FIXED\_PRECONDITIONER\_BINDING.md/json}.
원 \texttt{results/full/UNIFORM\_ARITHMETIC.json}의 $C,B,q,S$.

\bibitem{threshold_binding}
WU088\_HH PHYS15, \texttt{FT03\_THRESHOLD\_IDENTITY.json}.
Controlled fixture와 photon/HH threshold 세 leaves의 source graph 및 bits 동일성.

\bibitem{checks}
WU088\_HH PHYS15, \texttt{checks/}의 photon, sector, nonlinear-transfer 및
Wolfram exact algebra 기록. 시험용 manufactured 입력과 원 captured leaves의 역할을 각 실행에 명시했다.

\bibitem{review}
WU088\_HH PHYS15 독립 검토,
\texttt{review/DECISION.json}, \texttt{review/REVIEW\_KO.md},
\texttt{review/INDEPENDENT\_THEORY\_CHECKS.json}.
최종 판정 \texttt{PROMOTE\_SCOPED}, 2026-10-13 KST.
\end{thebibliography}
\endgroup


\end{document}

